\section{The Gap Between Standard KD Assumptions and Practice} \label{sec:12-introduction-gap}

\begin{foldable}
  The theoretical appeal of \ac{KD} rests on two assumptions that rarely hold outside carefully controlled research settings.
  The first is that the practitioner has access to the original training dataset, or at least a dataset drawn from the same distribution, so that meaningful queries to the teacher can be made during distillation.
  The second is that the teacher is a white-box system: its internal representations are fully observable, meaning the practitioner can read off the full output probability distribution over all classes, and optionally also the intermediate feature activations and gradients that pass through the network.
  Taken together, these assumptions frame \ac{KD} as a controlled lab procedure in which the distillation engineer has almost complete access to both data and model.
  In practice, neither assumption is routinely satisfied.
\end{foldable}

\begin{foldable}
  Consider first the data assumption.
  Training data for high-performing models is often the most commercially and legally sensitive asset held by an organization.
  Privacy legislation such as the General Data Protection Regulation (GDPR) in Europe and the Health Insurance Portability and Accountability Act (HIPAA) in the United States strictly limits the collection, retention, and redistribution of personal or medical records, making it legally impossible to share the data used to train a clinical imaging model with a third-party deployer.
  Even where regulation is not a binding constraint, intellectual property law and commercial confidentiality create strong incentives to withhold training corpora: a company that has spent years curating a proprietary dataset will not release it merely to facilitate compression of the model it trained.
  Beyond legal and strategic barriers, raw cost is a practical obstacle in its own right; large-scale datasets in computer vision and natural language processing can occupy terabytes of storage and require expensive curation pipelines, making them infeasible to distribute or even store locally at the deployment site.
  The net effect is that a practitioner wishing to distil a capable model frequently finds that the original training data has been destroyed, was never shared, or is simply too large to transfer.
\end{foldable}

\begin{foldable}
  The second assumption --- white-box teacher access --- is equally fragile.
  The dominant commercial model for deploying large neural networks has shifted decisively toward black-box inference \ac{API}{}s.
  GPT-style language models, vision classification services, and multimodal foundation models are made available as hosted endpoints: a client submits an input and receives a prediction, but the model weights, architecture, and internal activations never leave the server.
  Even the output is often severely truncated; many production \ac{API}{}s return only the top-predicted class label or a short ranked list, rather than the full probability distribution over thousands of categories.
  This hard-label regime is a deliberate design choice, since returning full softmax distributions would expose the model to higher risks of extraction attacks and reverse-engineering.
  The consequence for distillation is severe: without access to the teacher's soft probability vector, the student loses the rich inter-class similarity information --- what Hinton et al.\ referred to as ``dark knowledge'' --- that makes \ac{KD} substantially more effective than training on one-hot ground-truth labels alone.
\end{foldable}

\begin{foldable}
  The effect of these restrictions is not merely practical inconvenience --- it reshapes the statistical properties of whatever training signal remains available.
  \S\ref{sec:13-introduction-fidelity} examines why this matters and what it demands of any solution.
\end{foldable}
